% 1_intro.tex
\section{Introducción}
En este documento se abordan dos problemas clásicos de la mecánica celeste mediante métodos numéricos: el problema de $N$-cuerpos aplicado a la órbita del cometa Halley, y el problema restringido de tres cuerpos aplicado al sistema Tierra-Luna. 

En primer lugar, se simula la trayectoria del cometa Halley bajo la influencia gravitatoria del Sol, evaluando la conservación de la energía y el momento angular. Posteriormente, se introduce la perturbación gravitacional de Júpiter para cuantificar el fenómeno de la precesión orbital. 

En segundo lugar, se analiza el problema restringido de tres cuerpos, calculando el pozo de potencial efectivo del sistema Tierra-Luna para determinar las coordenadas de los cinco puntos de Lagrange ($L_1$ a $L_5$). Finalmente, se evalúa la estabilidad tridimensional de las órbitas alrededor de estos puntos de equilibrio frente a pequeñas perturbaciones cinemáticas. 